\documentclass[12pt]{article}
\usepackage[T2A,T1]{fontenc}
\usepackage[utf8]{inputenc}
\usepackage[english,ukrainian]{babel}
\usepackage[left=30mm,right=15mm,top=20mm,bottom=20mm]{geometry}
\usepackage{setspace}
\usepackage{parskip}
\usepackage{tikz}
\usetikzlibrary{svg.path}

\onehalfspacing
\setlength{\parindent}{1.25cm}

\newcommand{\tryzub}[1][1]{%
  \begin{tikzpicture}[scale=0.03*#1, yscale=-1, baseline=(current bounding box.center)]
    % Shield: transparent background, black stroke
    \filldraw[fill=white, draw=black, line width=1.5pt]
      svg {m5 5h650v689c0 48-29 97-76 117l-251 105-251-105c-44-20-76-65-72-117z};

    % Right half of Tryzub (Black)
    \fill[color=black] 
      svg {m329 53c-6 4-2 396 0 401 12 43 29 81 48 112-104 31-63 146-48 287 7-12 17-21 28-28 43-34 70-81 79-132h133v-580c-148 88-132 213-148 361 59-10 75 76-3 83-92-149-59-257-59-419 0-37-9-62-30-85zm200 143v297h-22c-6-23-21-41-42-51 8-85 10-189 64-246zm-22 337h22v120h-89c0-19-3-39-8-58 36-4 68-29 75-62zm-114 71c4 16 6 33 6 49h-50c1-27 19-44 44-49zm-44 89h46c-7 31-23 61-46 85z};

    % Left (mirrored) half of Tryzub (Black)
    \fill[color=black] 
      svg {m331 53c6 4 2 396 0 401-12 43-29 81-48 112 104 31 63 146 48 287-7-12-17-21-28-28-43-34-70-81-79-132h-133v-580c148 88 132 213 148 361-59-10-75 76 3 83 92-149 59-257 59-419 0-37 9-62 30-85zm-200 143v297h22c6-23 21-41 42-51-8-85-10-189-64-246zm22 337h-22v120h89c0-19 3-39 8-58-36-4-68-29-75-62zm114 71c-4 16-6 33-6 49h50c-1-27-19-44-44-49zm44 89h-46c7 31 23 61 46 85z};
  \end{tikzpicture}%
}


\begin{document}

\begin{center}
  \tryzub[0.8]\\[0.3cm]
  {\selectlanguage{ukrainian}\small\MakeUppercase{Державна судова адміністрація України}\\}
  {\selectlanguage{ukrainian}\small\bfseries\MakeUppercase{Державне підприємство «Інформаційні судові системи»}\\}
  \vspace{0.2cm}
  \hrule height 1.5pt
  \vspace{0.5cm}
  {\selectlanguage{ukrainian}\Large\bfseries НАКАЗ\\}
\end{center}

\vspace{0.2cm}

\noindent
\parbox[t]{0.33\textwidth}{\selectlanguage{ukrainian} 05.02.2024}%
\parbox[t]{0.33\textwidth}{\centering\selectlanguage{ukrainian} Київ}%
\parbox[t]{0.33\textwidth}{\raggedleft\selectlanguage{ukrainian} № 22\_\allowbreak{}ОД}

\vspace{0.8cm}

\noindent\textbf{Про організацію робіт із розробки та впровадження Комплексної системи захисту інформації}

\vspace{0.5cm}

З метою створення належних умов безпеки та захисту інформації в Єдиній судовій інформаційно-комунікаційній системі «ЄСІКС» (далі -- ЄСІКС), відповідно до Закону України «Про захист інформації в інформаційно-комунікаційних системах», Постанови Кабінету Міністрів України від 29.03.2006 № 373 «Про затвердження Правил забезпечення захисту інформації в інформаційних, телекомунікаційних та інформаційно-комунікаційних системах», вимог НД ТЗІ 3.7-001-99 «Методичні вказівки щодо розробки технічного завдання на створення комплексної системи захисту інформації в автоматизованій системі», та на виконання контролів планування (PL), оцінки ризиків (RA) та придбання послуг (SA),

\noindent\textbf{НАКАЗУЮ:}

1. Розпочати роботи з проектування та створення Комплексної системи захисту інформації (далі -- КСЗІ) в ЄСІКС.

2. Залучити в установленому законодавством порядку ліцензованого розробника ДП «УСС» для проведення обстеження середовища функціонування ЄСІКС, проведення аналізу ризиків та розробки Технічного завдання на створення КСЗІ.

3. Створити спільну робочу групу з розробки та впровадження КСЗІ у складі:
\begin{itemize}
  \item Керівник робочої групи -- Петро СИДОРЕНКО;
  \item Члени робочої групи -- представники СЗІ, IT-департаменту підприємства та представники розробника.
\end{itemize}

4. Робочій групі спільно з розробником ДП «УСС»:
\begin{itemize}
  \item у строк до 20.02.2024 провести інвентаризацію апаратних та програмних засобів ЄСІКС та скласти Акт обстеження середовища;
  \item у строк до 28.02.2024 провести первинну оцінку ризиків безпеки (RA-3) та визначити категорію безпеки інформації (RA-2);
  \item розробити та погодити в Держспецзв'язку Технічне завдання на створення КСЗІ (PL-2).
\end{itemize}

5. Фінансовому департаменту забезпечити фінансування зазначених робіт відповідно до умов укладеного договору з ДП «УСС».

6. Контроль за виконанням наказу покласти на керівника СЗІ.

\vspace{1.5cm}

\noindent
\parbox[t]{0.5\textwidth}{\selectlanguage{ukrainian}\bfseries Генеральний директор}%
\parbox[t]{0.5\textwidth}{\raggedleft\selectlanguage{ukrainian}\bfseries Василь ШЕВЧЕНКО}

\newpage

\begin{center}
  \textbf{\selectlanguage{ukrainian}АРКУШ ПОГОДЖЕННЯ}
\end{center}

\vspace{0.8cm}

\noindent
\textbf{ПОГОДЖЕНО:}

\vspace{0.5cm}
\noindent
\parbox[t]{0.5\textwidth}{\selectlanguage{ukrainian}Керівник робочої групи}%
\parbox[t]{0.5\textwidth}{\raggedleft\selectlanguage{ukrainian} Петро СИДОРЕНКО}

\vspace{0.8cm}
\noindent
\parbox[t]{0.5\textwidth}{\selectlanguage{ukrainian}Керівник СЗІ}%
\parbox[t]{0.5\textwidth}{\raggedleft\selectlanguage{ukrainian} \rule{4cm}{0.5pt}}

\vspace{0.8cm}
\noindent
\parbox[t]{0.5\textwidth}{\selectlanguage{ukrainian}Директор фінансового департаменту}%
\parbox[t]{0.5\textwidth}{\raggedleft\selectlanguage{ukrainian} \rule{4cm}{0.5pt}}

\end{document}
